% =====================================================================
\chapter{Modulo 5 -- Template di Funzione e di Classe}
% =====================================================================

% ---------------------------------------------------------------------
\section{Teoria \& Analogia Chiave}
% ---------------------------------------------------------------------

\subsection*{Il problema: lo stesso codice per tipi diversi}

Scrivi una funzione che restituisce il massimo tra due interi:
\begin{lstlisting}
int massimo(int a, int b) { return (a > b) ? a : b; }
\end{lstlisting}
Poi ti serve per i \cc|double|. Con l'overloading ne scrivi un'altra, identica
tranne che per il tipo. Poi per le \cc|string|, poi per i \cc|char|\ldots{} Hai
quattro copie dello stesso codice: se trovi un errore, va corretto quattro volte.
Peggio ancora con una classe: una pila di interi, una pila di stringhe, una pila
di rettangoli sarebbero tre classi quasi identiche.

Eppure l'algoritmo è lo stesso: ``se il primo è più grande restituisci il primo,
altrimenti il secondo''. L'unica cosa che cambia è il \textbf{tipo}. La
\textbf{programmazione generica} consiste nello scrivere algoritmi e strutture
dati \textbf{una volta sola}, lasciando il tipo come un parametro da specificare
dopo, quando il codice viene usato. In C++ lo strumento per farlo è il
\textbf{template} (``modello'', ``sagoma'').

\begin{analogia}[: lo stampo con il materiale da scegliere]
Un template è uno stampo per biscotti in cui \emph{non} hai ancora deciso
l'impasto. Lo stampo dice la forma; quando ti serve, scegli l'impasto (frolla,
cioccolato, pan di zenzero) e ottieni biscotti con quella forma e quel materiale.
Il template dice la ``forma'' del codice; il tipo concreto (\cc|int|,
\cc|string|, \cc|Rettangolo|) è l'impasto.
\end{analogia}

\subsection*{Template di funzione}

\begin{lstlisting}
template <typename T>          // "T è un tipo, che deciderò dopo"
T massimo(T a, T b) {
    return (a > b) ? a : b;
}
\end{lstlisting}
\begin{itemize}
  \item \cc|template <typename T>| va scritto \textbf{subito prima} della funzione.
        \cc|T| è un nome a scelta (per convenzione \cc|T|, ``Type''); al posto di
        \cc|typename| si può scrivere \cc|class|, con lo stesso significato.
  \item Dentro la funzione \cc|T| si usa come un tipo qualunque: per i parametri,
        per il valore di ritorno, per le variabili locali.
\end{itemize}

\textbf{Cosa fa il compilatore.} Un template \emph{non} è una funzione: è una
ricetta per generare funzioni. Quando scrivi \cc|massimo(3, 7)|, il compilatore
vede che gli argomenti sono \cc|int|, \textbf{deduce} \cc|T = int| e genera (si dice
\textbf{istanzia}) una vera funzione \cc|massimo<int>|, sostituendo \cc|int| a
ogni \cc|T|. Se poi scrivi \cc|massimo(2.5, 1.5)| genera anche
\cc|massimo<double>|. Tutto questo avviene \textbf{durante la compilazione}: i tipi
sono decisi prima che il programma giri, e i controlli sui tipi restano rigorosi.
Il codice sorgente contiene un solo template, ma il codice compilato contiene
\textbf{una copia per ogni tipo usato}.

\textbf{Quando la deduzione non basta.} \cc|massimo(3, 4.5)| non compila: dal
primo argomento \cc|T| sarebbe \cc|int|, dal secondo \cc|double|, e il
compilatore non sceglie da solo. In questi casi il tipo si indica esplicitamente
tra parentesi angolari: \cc|massimo<double>(3, 4.5)|. Il \cc|3| viene allora
convertito in \cc|3.0|.

\textbf{Il ``contratto'' implicito.} \cc|massimo| usa l'operatore \cc|>|. Quindi
funziona con qualunque tipo che \textbf{sappia fare} \cc|>|: \cc|int|, \cc|double|,
\cc|char|, \cc|string| (ordine alfabetico). Con una classe tua, ad esempio
\cc|Rettangolo|, il compilatore genererebbe codice che confronta due rettangoli
con \cc|>|: se la classe non definisce quell'operatore, ottieni un errore, spesso
lungo e confuso, che punta \emph{dentro} il template. Per far funzionare il
template devi ``insegnare'' alla classe l'operatore richiesto:
\begin{lstlisting}
class Rettangolo {
    double b, h;
public:
    double area() const { return b * h; }
    bool operator<(const Rettangolo& altro) const {   // "r1 < r2" = confronto le aree
        return area() < altro.area();
    }
};
\end{lstlisting}
È un'anteprima del Modulo 10: per ora basta sapere che \cc|operator<| è un metodo
che il compilatore chiama quando scrivi \cc|r1 < r2|. Allo stesso modo, per
stampare un oggetto con \cc|cout <<| serve \cc|operator<<|; negli esempi di questo
modulo lo trovi già scritto, come formula da copiare. Un buon consiglio: nei tuoi
template usa il \textbf{minor numero possibile} di operatori (ad esempio solo
\cc|<|), così i tipi dovranno soddisfare meno requisiti.

\textbf{Più parametri e parametri non di tipo.} Un template può avere più tipi
(\cc|template <typename T, typename U>|) e anche parametri che sono
\textbf{valori} costanti, ad esempio un intero:
\cc|template <typename T, int N>|. Si possono dare valori di default:
\cc|template <typename T = int, int N = 5>|.

\textbf{Specializzazione.} A volte la versione generale è sbagliata per un tipo
particolare. Confrontare due \cc|double| con \cc|==| è inaffidabile
(\cc|0.1 + 0.2 == 0.3| è falso!); confrontare due \cc|const char*| con \cc|==|
confronta gli indirizzi, non il testo. Si può allora scrivere una versione
apposita per quel tipo, che ha la precedenza su quella generale:
\begin{lstlisting}
template <typename T>
bool uguali(T a, T b) { return a == b; }          // versione generale

template <>
bool uguali<double>(double a, double b) {         // SOLO per double
    double d = a - b;
    return (d < 0 ? -d : d) < 1e-9;
}
\end{lstlisting}

\subsection*{Template di classe}

Lo stesso meccanismo vale per le classi. È il caso più importante per l'esame:
liste, pile, code e alberi dei moduli 11--13 si scrivono come template.
\begin{lstlisting}
template <typename T>
class Contenitore {
    T dati[50];
    int elementi;
public:
    Contenitore() : elementi(0) {}
    void inserire(T el);
    T estrarre();
};
Contenitore<int> interi;        // il tipo va SEMPRE scritto tra < >
Contenitore<float> reali;       // per le classi non c'è deduzione
\end{lstlisting}
\cc|Contenitore| da solo non è un tipo: è un modello. \cc|Contenitore<int>| e
\cc|Contenitore<float>| sono due \textbf{classi diverse}, generate dal
compilatore, che non hanno nulla in comune (non puoi assegnare l'una all'altra).

\textbf{Metodi definiti fuori dalla classe.} Ogni metodo scritto fuori va
preceduto di nuovo da \cc|template <typename T>|, e il nome della classe va scritto
con il parametro:
\begin{lstlisting}
template <typename T>
void Contenitore<T>::inserire(T el) {
    dati[elementi++] = el;
}
\end{lstlisting}
Dentro la classe (e dentro i suoi metodi) si può scrivere semplicemente
\cc|Contenitore| invece di \cc|Contenitore<T>|.

\textbf{Template e memoria dinamica.} Un template di classe che usa \cc|new T[n]|
segue esattamente le regole del Modulo 4: distruttore, costruttore di copia,
\cc|operator=|. In più, \cc|new T[n]| richiede che \cc|T| abbia un costruttore di
default, perché crea subito \cc|n| oggetti.

\subsection*{Dove si scrivono i template: tutto nell'header}

Nel Modulo 2 abbiamo diviso le classi in \texttt{.h} (dichiarazione) e \texttt{.cpp}
(implementazione). Con i template \textbf{non si può fare}. Il motivo segue da come
lavorano: il compilatore genera \cc|Contenitore<int>::inserire| solo quando, nello
stesso file che sta compilando, vede sia l'uso (\cc|Contenitore<int>|) sia il corpo
del template. Se il corpo sta in \texttt{contenitore.cpp} e l'uso in
\texttt{main.cpp}, quando compila \texttt{main.cpp} conosce solo la dichiarazione;
quando compila \texttt{contenitore.cpp} non sa che serve la versione \cc|int|.
Nessuno dei due genera il codice, e il linker si ferma con \emph{undefined reference
to `Contenitore<int>::inserire(int)'}.

La soluzione standard: \textbf{tutto il template (dichiarazione e corpi dei metodi)
va nel file header}, che viene incluso dove serve.

\begin{inparole}
\begin{itemize}
  \item Un template è una ricetta: \cc|template <typename T>| e poi codice in cui
        \cc|T| è un tipo qualunque.
  \item Il compilatore genera una versione per ogni tipo usato (istanziazione), a
        compile-time.
  \item Per le funzioni il tipo di solito si deduce; altrimenti
        \cc|f<double>(...)|. Per le classi si scrive sempre: \cc|Pila<int>|.
  \item Il tipo deve supportare tutte le operazioni usate nel template (\cc|<|,
        \cc|==|, \cc|<<|, \ldots); per le tue classi le devi definire tu.
  \item Fuori dalla classe: \cc|template <typename T>| + \cc|Classe<T>::metodo|.
  \item I template vanno interamente nell'header, non in un \texttt{.cpp} separato.
\end{itemize}
\end{inparole}

% ---------------------------------------------------------------------
\section{Rappresentazione Visiva della Memoria (TikZ)}
% ---------------------------------------------------------------------

\textbf{Un template, tante funzioni.} Il compilatore genera una versione per ogni
tipo con cui il template viene usato; nell'eseguibile finiscono solo quelle.
\begin{center}
\begin{tikzpicture}
  \node[processo,text width=4.2cm] (t) at (0,0) {\texttt{template <typename T>}\\\texttt{T massimo(T a, T b)}};
  \node[etichetta,above=0.1cm of t] {codice sorgente (uno solo)};
  \node[file,text width=5.4cm,font=\footnotesize\ttfamily,fill=green!15!sfondo] (i) at (9.2,1.6) {int massimo(int, int)};
  \node[file,text width=5.4cm,font=\footnotesize\ttfamily,fill=green!15!sfondo] (d) at (9.2,0) {double massimo(double, double)};
  \node[file,text width=5.4cm,font=\footnotesize\ttfamily,fill=green!15!sfondo] (s) at (9.2,-1.6) {string massimo(string, string)};
  \node[etichetta,above=0.1cm of i] {codice compilato (una copia per tipo)};
  \draw[freccia] (t.east) -- (i.west);
  \draw[freccia] (t.east) -- (d.west);
  \draw[freccia] (t.east) -- (s.west);
  \node[etichetta,anchor=south] at ($(i.west)+(-2.0,0.12)$) {\texttt{massimo(3, 7)}};
  \node[etichetta,anchor=south] at ($(d.west)+(-2.0,0.12)$) {\texttt{massimo(2.5, 1.5)}};
  \node[etichetta,anchor=north] at ($(s.west)+(-2.0,-0.12)$) {\texttt{massimo(s1, s2)}};
\end{tikzpicture}
\end{center}

\needspace{6cm}
\textbf{Un template di classe, classi indipendenti.} \cc|Scatola<int>| e
\cc|Scatola<string>| del pattern hanno la stessa struttura ma sono due tipi diversi,
con array di elementi di tipo diverso nello heap.
\begin{center}
\begin{tikzpicture}
  \node[cella,fill=stackcol,minimum width=2.6cm] (si) at (0,0.8) {si: Scatola<int>};
  \node[cella,fill=stackcol,minimum width=2.6cm] (ss) at (0,-0.8) {ss: Scatola<string>};
  \foreach \v [count=\i] in {1,4,9,16,25,36,{},{}}
    \node[cellina,fill=heapcol] (a\i) at (3.2+\i*0.75,0.8) {\v};
  \node[cella,fill=heapcol,minimum width=1.4cm] (b1) at (4.6,-0.8) {\q{alfa}};
  \node[cella,fill=heapcol,minimum width=1.4cm,right=0cm of b1] (b2) {\q{beta}};
  \node[cella,fill=heapcol,minimum width=1cm,right=0cm of b2] (b3) {};
  \node[cella,fill=heapcol,minimum width=1cm,right=0cm of b3] (b4) {};
  \draw[freccia] (si.east) -- (a1.west);
  \draw[freccia] (ss.east) -- (b1.west);
  \node[etichetta,right=0.1cm of a8] {8 \texttt{int}};
  \node[etichetta,right=0.1cm of b4] {4 \texttt{string}};
\end{tikzpicture}
\end{center}

\needspace{6cm}
\textbf{Perché un template non si separa in \texttt{.h} e \texttt{.cpp}.}
\begin{center}
\begin{tikzpicture}
  \node[file,text width=3.4cm] (h) at (0,1.2) {pila.h\\ \footnotesize dichiarazione};
  \node[file,text width=3.4cm] (c) at (0,-0.6) {pila.cpp\\ \footnotesize corpi dei metodi};
  \node[file,text width=3.4cm] (m) at (0,-2.2) {main.cpp\\ \footnotesize usa \texttt{Pila<int>}};
  \node[processo,text width=5.6cm] (cc) at (6.6,-0.6) {compilando \texttt{pila.cpp}: \emph{nessuno usa Pila<int>, non genero nulla}};
  \node[processo,text width=5.6cm] (cm) at (6.6,-2.2) {compilando \texttt{main.cpp}: \emph{non vedo i corpi, non posso generare}};
  \draw[freccia] (c) -- (cc);
  \draw[freccia] (m) -- (cm);
  \node[etichetta,text=red!75!black,text width=8cm,align=center] at (6.6,-3.4) {linker: \emph{undefined reference to Pila<int>::push(int)}};
\end{tikzpicture}
\end{center}

% ---------------------------------------------------------------------
\section{Pattern di Codice di Laboratorio}
% ---------------------------------------------------------------------

Template di funzione (\cc|massimo|, \cc|scambia|, \cc|stampaArray|) e un template di
classe con memoria dinamica e Regola dei Tre, con i metodi definiti fuori dalla
classe: è la forma esatta che userai per liste e pile.

\codice{codice/m05/pattern.cpp}
\uscita{codice/m05/output/pattern.txt}

\textbf{Cosa osservare:}
\begin{itemize}
  \item \cc|massimo(string("ciao"), string("mondo"))|: se scrivessi
        \cc|massimo("ciao", "mondo")|, \cc|T| diventerebbe \cc|const char*| e il
        confronto avverrebbe tra \textbf{indirizzi} (vedi l'Esercizio~5.9).
  \item \cc|scambia(a, b)| e \cc|scambia(s1, s2)| generano due funzioni diverse,
        \cc|scambia<int>| e \cc|scambia<string>|.
  \item Ogni metodo di \cc|Scatola| scritto fuori dalla classe ripete
        \cc|template <typename T>| e usa \cc|Scatola<T>::|.
\end{itemize}

% ---------------------------------------------------------------------
\section{Tips \& Trappole da Esame}
% ---------------------------------------------------------------------

\begin{trappola}[: template diviso in .h e .cpp]
Mettere i corpi dei metodi di un template in un \texttt{.cpp} separato produce
\emph{undefined reference} al linker. Tutto il template va nell'header (o tutto
nello stesso file del \cc|main|, come negli esercizi d'esame su file unico).
\end{trappola}

\begin{trappola}[: dimenticare template <typename T> fuori dalla classe]
Ogni metodo definito fuori deve iniziare con \cc|template <typename T>| e usare
\cc|NomeClasse<T>::|. Scrivere solo \cc|NomeClasse::metodo| dà errori come
\emph{'NomeClasse' is not a class, namespace, or enumeration}.
\end{trappola}

\begin{trappola}[: deduzione ambigua]
\cc|massimo(3, 4.5)| non compila (\emph{deduced conflicting types for parameter 'T'}).
Usa \cc|massimo<double>(3, 4.5)| oppure due parametri di tipo diversi.
\end{trappola}

\begin{trappola}[: la classe non supporta l'operatore]
Un template che usa \cc|<| applicato a una tua classe senza \cc|operator<| dà un
errore che indica la riga \emph{dentro il template} (\emph{no match for 'operator<'}).
Il problema però è nella classe: aggiungi l'operatore.
\end{trappola}

\begin{trappola}[: stringhe C nei template]
Con \cc|"ciao"| \cc|T| è \cc|const char*|: \cc|==| e \cc|<| confrontano indirizzi.
Usa \cc|string| o scrivi una specializzazione.
\end{trappola}

\begin{consiglio}[: scrivere prima la versione con un tipo concreto]
Se un template non ti viene, scrivi prima la funzione o la classe per \cc|int|,
provala, poi sostituisci \cc|int| con \cc|T| dove serve e aggiungi
\cc|template <typename T>|. Gli errori di compilazione dei template sono lunghi:
leggi la \textbf{prima} riga che contiene il tuo file e il tuo numero di riga.
\end{consiglio}

% ---------------------------------------------------------------------
\section{Quiz a Crocette (Prova Scritta)}
% ---------------------------------------------------------------------

\begin{quiz}{Quando vengono generate le versioni concrete di un template?}
  \item durante l'esecuzione, la prima volta che vengono chiamate
  \item durante la compilazione, per ogni tipo con cui il template viene usato
  \item durante il linking, per tutti i tipi possibili
  \item mai: il template viene interpretato
\end{quiz}

\begin{quiz}{Dato \texttt{template <typename T> T massimo(T a, T b);} quale chiamata \textbf{non} compila?}
  \item \texttt{massimo(3, 7)}
  \item \texttt{massimo(2.5, 1.0)}
  \item \texttt{massimo(3, 4.5)}
  \item \texttt{massimo<double>(3, 4.5)}
\end{quiz}

\begin{quiz}{In \texttt{template <typename T>}, la parola \texttt{typename} può essere sostituita da:}
  \item \texttt{type}
  \item \texttt{class}
  \item \texttt{struct}
  \item \texttt{auto}
\end{quiz}

\begin{quiz}{Come si dichiara un oggetto della classe template \texttt{Pila} per contenere stringhe?}
  \item \texttt{Pila p<string>;}
  \item \texttt{Pila<string> p;}
  \item \texttt{Pila p(string);}
  \item \texttt{template Pila<string> p;}
\end{quiz}

\begin{quiz}{Un template \texttt{ordina(T a[], int n)} usa l'operatore \texttt{<}. Cosa serve per ordinare un array di \texttt{Rettangolo}?}
  \item niente: il compilatore confronta i rettangoli byte per byte
  \item che \texttt{Rettangolo} definisca \texttt{operator<}
  \item una specializzazione obbligatoria \texttt{ordina<Rettangolo>}
  \item che \texttt{Rettangolo} sia una \texttt{struct}
\end{quiz}

\begin{quiz}{Come si definisce fuori dalla classe il metodo \texttt{push} della classe template \texttt{Pila<T>}?}
  \item \texttt{void Pila::push(T x) \{...\}}
  \item \texttt{template <typename T> void Pila<T>::push(T x) \{...\}}
  \item \texttt{void Pila<T>::push(T x) \{...\}}
  \item \texttt{template void Pila::push(T x) \{...\}}
\end{quiz}

\begin{quiz}{Un template di classe è dichiarato in \texttt{pila.h} e i suoi metodi sono definiti in \texttt{pila.cpp}. Il \texttt{main} usa \texttt{Pila<int>}. Cosa succede tipicamente?}
  \item funziona tutto normalmente
  \item errore del linker: \emph{undefined reference}
  \item errore del preprocessore
  \item il programma funziona ma è più lento
\end{quiz}

\begin{quiz}{\texttt{Contenitore<int>} e \texttt{Contenitore<double>} sono:}
  \item la stessa classe con due nomi
  \item due classi distinte generate dallo stesso template
  \item due oggetti della stessa classe
  \item una classe e una sua sottoclasse
\end{quiz}

\begin{quiz}{Cosa indica \texttt{template <>} prima di \texttt{bool uguali<double>(double a, double b)}?}
  \item un template senza parametri, quindi un errore
  \item una specializzazione esplicita del template per il tipo \texttt{double}
  \item un template che accetta qualunque tipo tranne \texttt{double}
  \item una funzione template con parametri di default
\end{quiz}

\begin{quiz}{In \texttt{template <typename T, int N> class Buffer \{ T dati[N]; \};} cos'è \texttt{N}?}
  \item un secondo tipo
  \item un parametro non di tipo: un valore intero costante noto a compile-time
  \item una variabile che si può cambiare durante l'esecuzione
  \item un errore: i template accettano solo tipi
\end{quiz}

\begin{soluzione}[: risposte commentate]
\begin{enumerate}[label=\textbf{\arabic*.}]
  \item \textbf{B}. L'istanziazione avviene a compile-time, solo per i tipi effettivamente usati.
  \item \textbf{C}. \cc|T| dedotto in modo contraddittorio (\cc|int| e \cc|double|). La D
        risolve indicandolo esplicitamente.
  \item \textbf{B}. \cc|template <class T>| è equivalente.
  \item \textbf{B}. Per le classi il tipo si specifica sempre tra \cc|< >| dopo il nome.
  \item \textbf{B}. Il template funziona con qualunque tipo che fornisca le operazioni
        usate. La specializzazione (C) è possibile ma non obbligatoria.
  \item \textbf{B}. Servono sia \cc|template <typename T>| sia \cc|Pila<T>::|.
  \item \textbf{B}. Nessuna unità di traduzione genera il codice di \cc|Pila<int>|.
  \item \textbf{B}. Ogni istanziazione è una classe a sé.
  \item \textbf{B}. È la sintassi della specializzazione esplicita.
  \item \textbf{B}. Infatti si può usare come dimensione di un array statico.
\end{enumerate}
\end{soluzione}

% ---------------------------------------------------------------------
\section{Laboratorio a Step (Simulazione d'Esame)}
% ---------------------------------------------------------------------

\begin{traccia}[: magazzino generico]
Si vuole scrivere un piccolo insieme di strumenti generici (funzioni e classi
template) per memorizzare, ordinare e stampare oggetti di qualunque tipo, e usarli
con interi, caratteri, stringhe e una classe \cc|Rettangolo|. Niente STL.
\end{traccia}

\begin{step}{template di funzione}
Scrivi \cc|template <typename T> void stampa(const T v[], int n)| (elementi separati
da virgole) e \cc|int indiceMassimo(const T v[], int n)|, che usa \textbf{solo}
l'operatore \cc|<|. Provale con un array di \cc|int| e uno di \cc|string|.\\[0.3em]
\textbf{Concetto nuovo:} \cc|template <typename T>|, deduzione del tipo.
\end{step}

\begin{step}{ordinamento generico}
Scrivi \cc|void ordina(T a[], int n)| (selection sort, usando solo \cc|<|) e
applicalo ad array di \cc|int|, \cc|char| e \cc|string|.\\[0.3em]
\textbf{Concetto nuovo:} un algoritmo, molti tipi.
\end{step}

\begin{step}{una classe compatibile}
Scrivi \cc|Rettangolo| (base, altezza, \cc|area()|) con \cc|operator<| che confronta
le aree e \cc|operator<<| per la stampa (usa la forma fornita negli esempi). Ordina
un array di rettangoli e trova il più grande con le funzioni degli step 1--2.
(È l'esercizio delle slide del corso.)\\[0.3em]
\textbf{Concetto nuovo:} requisiti che un tipo deve soddisfare per un template.
\end{step}

\begin{step}{template di classe}
Scrivi \cc|template <typename T> class Magazzino| con un array dinamico che
raddoppia, Regola dei Tre, \cc|inserire|, \cc|estrarre| (lancia
\cc|underflow_error| se vuoto), \cc|numero|, \cc|vuoto|, \cc|ordinaTutto()| (riusa
\cc|ordina|) e \cc|stampaTutto()|. Usalo con \cc|double| e con \cc|Rettangolo|.\\[0.3em]
\textbf{Concetto nuovo:} template di classe con memoria dinamica.
\end{step}

\begin{step}{due parametri e template annidati}
Scrivi \cc|template <typename K, typename V> class Coppia| (chiave e valore,
\cc|operator<| sulla chiave, \cc|operator<<|) e crea un
\cc|Magazzino< Coppia<string, int> >| come inventario: inserisci tre articoli,
ordinali per nome e stampali.\\[0.3em]
\textbf{Concetto nuovo:} più parametri di tipo, istanziazione con un tipo template.
\end{step}

\subsection*{Soluzione completa}
\codice{codice/m05/lab.cpp}
\uscita{codice/m05/output/lab.txt}

\begin{consiglio}[: cosa notare nella soluzione]
\begin{itemize}
  \item Nello step 3 \cc|6x0.5| (area 3) viene prima di \cc|2x5| (area 10): l'ordine
        è per area, non per base.
  \item \cc|Coppia()| inizializza con \cc|chiave()| e \cc|valore()|: per \cc|string|
        è la stringa vuota, per \cc|int| è 0. Serve perché \cc|new T[cap]| crea
        subito oggetti con il costruttore di default.
  \item \cc|Magazzino< Coppia<string, int> >|: lo spazio tra le due \cc|>| finali
        era obbligatorio nel vecchio C++98 (altrimenti \cc|>>| era letto come
        operatore); oggi non serve, ma non fa danni.
\end{itemize}
\end{consiglio}

% ---------------------------------------------------------------------
\section{Domande e Risposte Modello per l'Orale}
% ---------------------------------------------------------------------

\begin{orale}[D1. Cos'è un template e a cosa serve?]
Un template è un modello di funzione o di classe parametrizzato su uno o più tipi,
che permette di scrivere un algoritmo o una struttura dati una sola volta,
indipendentemente dal tipo degli elementi su cui opera: è lo strumento della
programmazione generica in C++. Il tipo viene specificato come se fosse un
parametro. Si evita così di duplicare codice identico per tipi diversi, mantenendo
la tipizzazione statica: i tipi sono determinati e controllati a compile-time.
\end{orale}

\begin{orale}[D2. Cosa significa istanziare un template?]
Significa che il compilatore, quando incontra un uso del template con un tipo
concreto, genera una versione specializzata del codice sostituendo il parametro con
quel tipo. Per i template di funzione il tipo viene di solito dedotto dagli
argomenti, ma si può indicare esplicitamente con \cc|f<tipo>(...)|; per i template
di classe va sempre indicato (\cc|Pila<int>|). Il codice sorgente contiene un solo
template, quello compilato contiene una copia per ogni tipo usato.
\end{orale}

\begin{orale}[D3. Perché i template si scrivono interamente negli header?]
Perché per istanziare un template il compilatore deve vederne la definizione
completa nel punto in cui viene usato. Se i corpi dei metodi stessero in un
\texttt{.cpp} separato, compilando il file che usa il template il compilatore
vedrebbe solo le dichiarazioni, e compilando il \texttt{.cpp} del template non
saprebbe per quali tipi generarlo: il risultato è un errore del linker. Per questo
dichiarazione e definizione si mettono insieme nell'header.
\end{orale}

\begin{orale}[D4. Cos'è la specializzazione di un template?]
È una versione alternativa del template scritta per un tipo particolare, che il
compilatore usa al posto di quella generale quando il parametro è proprio quel
tipo. Si scrive con \cc|template <>| seguito dalla funzione o classe con il tipo
indicato. Si usa quando l'algoritmo generale non è corretto o efficiente per quel
tipo: ad esempio il confronto con \cc|==| tra \cc|double|, che va fatto con una
tolleranza, o tra stringhe C, dove \cc|==| confronterebbe gli indirizzi.
\end{orale}

\begin{orale}[D5. Che requisiti deve avere un tipo per essere usato in un template?]
Deve supportare tutte le operazioni che il template usa sugli oggetti di quel tipo:
operatori di confronto, aritmetici, di stampa, costruttore di default se il template
crea array con \cc|new T[n]|, costruttore di copia e assegnamento se li copia. Se un
requisito manca, l'errore emerge in compilazione al momento dell'istanziazione. Per
le classi definite dall'utente bisogna quindi sovraccaricare gli operatori necessari.
\end{orale}

% ---------------------------------------------------------------------
\section{Esercizi al PC}
% ---------------------------------------------------------------------

\begin{esercizio}{Massimo e minimo}{facile}
Scrivi i template \cc|massimo| e \cc|minimo| e provali con \cc|int|, \cc|double|,
\cc|char| e \cc|string|. Prova a chiamare \cc|massimo(3, 4.5)|, leggi l'errore, poi
risolvi in due modi diversi.
\end{esercizio}

\begin{esercizio}{Scambia e ruota}{facile}
Scrivi \cc|template <typename T> void scambia(T&, T&)| e, usandola,
\cc|void ruota(T& a, T& b, T& c)| che sposta i valori a sinistra (\cc|a| prende il
valore di \cc|b|, \cc|b| quello di \cc|c|, \cc|c| quello di \cc|a|). Provala con
interi, stringhe e reali.
\end{esercizio}

\begin{esercizio}{Stampa e somma}{facile}
Scrivi \cc|stampaArray| e \cc|T somma(const T v[], int n)|, partendo da
\cc|T s = T();|. Provala con \cc|int|, \cc|double| e \cc|string|: cosa vuol dire
``sommare'' delle stringhe?
\end{esercizio}

\begin{esercizio}{Contare le occorrenze}{facile}
Scrivi \cc|int conta(const T v[], int n, const T& x)| e un template con due
parametri di tipo, \cc|int contaMaggiori(const T v[], int n, U soglia)|, che
permette di confrontare un array di \cc|double| con una soglia \cc|int|.
\end{esercizio}

\begin{esercizio}{Ricerca generica}{medio}
Scrivi \cc|int ricerca(T dato, T v[], int dim)| (indice della prima occorrenza, o
$-1$) e una classe \cc|Punto| con \cc|operator==|. Cerca nei tre array (\cc|int|,
\cc|char|, \cc|Punto|) valori letti da tastiera. (Esercizio delle slide.)
\end{esercizio}

\begin{esercizio}{Ordinamento generico}{medio}
Scrivi \cc|void ordina(T a[], int n)| (bubble sort, usando solo \cc|<|) e ordinaci
un array di \cc|int|, uno di \cc|char| e uno di \cc|Rettangolo| confrontati per area.
Stampa prima e dopo. (Esercizio delle slide.)
\end{esercizio}

\begin{esercizio}{Coppia di due tipi}{medio}
Scrivi \cc|template <typename T, typename U> class Coppia| con getter, setter,
\cc|stampa()| e un metodo \cc|invertita()| che restituisce una \cc|Coppia<U, T>|.
Scrivi anche una funzione template \cc|creaCoppia(a, b)| che deduce i tipi da sola.
\end{esercizio}

\begin{esercizio}{Contenitore a dimensione fissa}{medio}
Scrivi \cc|template <typename T = int, int N = 5> class Contenitore| con un array
\textbf{statico} \cc|T dati[N]|, \cc|inserire| (lancia \cc|overflow_error| se pieno),
\cc|estrarre| (lancia \cc|underflow_error| se vuoto), \cc|numero|, \cc|vuoto|,
\cc|capacita|. Crea un \cc|Contenitore<>| e un \cc|Contenitore<char, 3>|.
\end{esercizio}

\begin{esercizio}{Specializzazioni}{difficile}
Scrivi \cc|template <typename T> bool uguali(T a, T b)| con \cc|==| e due
specializzazioni: per \cc|double| (tolleranza $10^{-9}$) e per \cc|const char*|
(confronto carattere per carattere). Mostra che \cc|0.1 + 0.2 == 0.3| è falso ma
\cc|uguali(0.1 + 0.2, 0.3)| è vero.
\end{esercizio}

\begin{esercizio}{Pila generica}{difficile}
Scrivi \cc|template <typename T> class Pila| su array dinamico (Regola dei Tre,
\cc|push|, \cc|pop| con eccezione, \cc|vuota|). Usala in una funzione template
\cc|invertiConPila(T v[], int n)| e per controllare se le parentesi di
un'espressione come \texttt{\{[(a+b)*c]-(d)\}} sono bilanciate.
\end{esercizio}

\subsection*{Soluzioni}
Le soluzioni commentate, con il codice e l'output reale, sono nel PDF separato
\textbf{Soluzioni} (\texttt{soluzioni-light.pdf} / \texttt{soluzioni-dark.pdf}),
allo stesso numero di esercizio. Il codice è anche in \texttt{codice/m05/}.
